\chapter*{ABSTRACT}
\addcontentsline{toc}{chapter}{ABSTRACT}

This degree project proposal aims to ground and structure the design of a technical support management and traceability computer system powered by artificial intelligence, tailored for computer service providers in the city of Tarija during the 2026 term. The research stems from the need to analyze how technical support service providers in the city of Tarija manage equipment intake, diagnostic records, budget communication, and repair tracking. Through field research, the study will identify the recording methods currently used and potential communication and information access difficulties experienced by clients, in order to ground the requirements of the proposed system. The methodology adopts a mixed-methods approach (qualitative and quantitative), with a descriptive-propositive level and a non-experimental cross-sectional design, utilizing a non-probabilistic convenience sampling strategy across two analysis units: 10 to 15 workshop technicians or owners and 50 to 100 computer service clients. The proposed software engineering solution adopts a Model-View-Controller (MVC) architecture utilizing the PHP Laravel framework, MySQL, and a responsive Bootstrap 5 interface, where the design contemplates a traceability module featuring sequential service statuses, photographic evidence, and soft-delete database constraints. Furthermore, it proposes integrating a large language model (LLM) conversational assistant through an external provider REST API, taking Google Gemini as a reference technological alternative, with the final model selection subject to service availability, continuity, usage costs, and Spanish processing capability. The Laravel backend will validate requests, query strictly authorized data, and provide the assistant with the necessary context to handle frequent inquiries and status lookups as a designed architectural capability. The proposal is grounded in IT governance frameworks such as ITIL and international standards including ISO/IEC 25010 and ISO/IEC 27001.
